\documentclass[UTF8,11pt,a4paper,fontset=none]{ctexart}

\setCJKmainfont[AutoFakeBold=2.5]{Songti SC}
\setCJKsansfont[AutoFakeBold=2.5]{PingFang SC}
\setCJKmonofont{Hiragino Sans GB}

\usepackage[margin=25mm]{geometry}
\usepackage{amsmath,amssymb,mathtools,bm}
\usepackage{booktabs,tabularx,array,multirow}
\usepackage{graphicx}
\usepackage{float}
\usepackage{xcolor}
\usepackage{enumitem}
\usepackage[colorlinks=true,linkcolor=black,citecolor=black,urlcolor=blue]{hyperref}

\newcolumntype{Y}{>{\raggedright\arraybackslash}X}
\newcommand{\method}{NCMA}
\newcommand{\clip}{CLIP}
\newcommand{\R}{\mathbb{R}}
\newcommand{\M}{\mathcal{M}}
\newcommand{\norm}[1]{\left\lVert #1 \right\rVert}
\newcommand{\softmax}{\operatorname{softmax}}
\newcommand{\topk}{\operatorname{topk}}
\newcommand{\Norm}{\operatorname{Norm}}

\setlist[itemize]{leftmargin=2em,itemsep=2pt,topsep=3pt}
\setlist[enumerate]{leftmargin=2em,itemsep=2pt,topsep=3pt}

\title{\bfseries 从异常语义到正常性边界：\\CLIP 零样本异常定位的实例条件化建模}
\author{匿名作者}
\date{}

\begin{document}

\maketitle

\begin{abstract}
\clip{} 零样本异常检测通常被表述为 normal 与 abnormal 状态之间的语义判别：给定未见目标类别，模型需要利用可迁移的状态文本完成图像级异常判别与像素级缺陷定位。
现有 CLIP-ZSAD 方法已经推进了这一语义迁移过程，通过状态提示、局部视觉-文本对齐、视觉上下文或轻量适配器，使 abnormal concept 能够脱离具体对象类别而被识别。
然而，异常定位并不只是异常语义识别。
在工业图像中，缺陷通常表现为局部外观相对于当前实例正常状态的越界；同一类颜色偏移、纹理扰动或几何变化，在一个实例中可能是合法正常变化，在另一个实例中却可能构成缺陷。
因此，异常语义能够跨类别迁移，并不意味着正常性边界也能由同一组 normal/abnormal 文本锚直接迁移到每个目标实例。
本文将 CLIP 零样本异常定位重构为实例条件化正常性边界估计，并提出正常变化流形适应（Normal-Change Manifold Adaptation, \method{}）。
\method{} 从当前图像中选择语义正常且空间一致的 patch，构建由局部 anchor、可容许变化方向和残差尺度组成的正常变化流形。
该流形刻画当前实例中 normal 语义的可接受范围，并将每个 patch 分解为边界内变化与边界外残差：前者用于估计当前实例的 normality boundary，后者作为异常定位证据保留。
在 MVTec AD 上，\method{} 将 AnomalyCLIP 的平均 P-AUROC/P-AUPRO/I-AUROC/I-AP 从 91.1/81.4/91.6/96.4 提升到 92.9/84.6/92.4/96.7。
这些结果说明，在不使用目标训练图像、标签、mask 或目标正常样本库的条件下，实例条件化正常性边界能够补足 CLIP-ZSAD 中固定状态比较所缺少的边界信息。
\end{abstract}

\noindent\textbf{关键词：} 零样本异常检测；视觉语言模型；实例条件化正常性边界；正常变化流形；异常定位

\section{引言}
\label{sec:introduction}

工业异常检测的核心不是识别一个新的对象类别，而是判断局部外观是否仍处在正常生产状态允许的范围内。
划痕、污染、结构缺失和部件错位往往只改变局部区域，却会使该区域越过正常性边界。
零样本异常检测进一步要求模型在目标类别没有训练图像和人工标注的条件下完成图像级异常判别与像素级定位。
这一设置对应新产线快速部署、缺陷形态开放和标注成本高昂等实际场景，也使预训练视觉语言模型成为自然的基础模型。

\clip{} 通过大规模图文对比学习获得开放词汇识别能力 \cite{radford2021clip}。
近期 CLIP-ZSAD 方法将这种能力转化为 normal/abnormal 状态建模，代表性路线包括状态词模板、窗口级局部证据、对象无关提示、动态提示、视觉上下文提示和轻量适配器 \cite{jeong2023winclip,chen2023aprilgan,zhou2024anomalyclip,yao2024adaclip,qu2024vcpclip,ma2025aaclip}。
这些工作推动了异常语义迁移：模型可以在未见目标类别上获得对象无关的 abnormal 表达，并把局部视觉证据投射到 normal/abnormal 状态空间中。
从这一角度看，CLIP-ZSAD 已经把问题推进到一个更细的层面：异常语义可以迁移，但异常定位仍要决定 normal 语义在当前实例中覆盖到哪里。

异常语义的可迁移性并不自动带来异常定位所需的边界判定。
异常定位面对的是另一个问题：当前图像中的局部偏离，是否仍能被该实例的正常状态接受。
相同的纹理方向变化、材质色差、局部几何扰动或成像波动，在不同实例中具有不同的正常含义。
如果模型始终使用同一组 normal/abnormal 语义锚评估所有目标图像，就隐含地假设正常性边界可以随语义一起固定迁移。
这一假设与工业图像的实例差异不匹配，因为正常性边界往往由当前实例内部的可接受变化决定，而不是由全局文本锚单独决定。

本文因此关注 CLIP-ZSAD 中被固定状态比较遮蔽的边界迁移问题。
我们将零样本异常定位重构为实例条件化正常性边界估计：CLIP 提供可迁移的状态语义，当前图像中的可靠正常证据定义该语义在实例上的可接受范围。
在这个视角下，normal 不再只是一个文本状态或单个视觉原型，而是一个局部边界结构。
它需要同时回答三个问题：正常外观的中心在哪里，哪些局部变化方向仍被允许，以及偏离到何种尺度后应被视为异常证据。

基于这一认识，本文提出正常变化流形适应（\method{}）。
\method{} 使用冻结的 \clip{} normal/abnormal 语义锚在当前图像中筛选语义正常且空间一致的候选 patch，并将这些候选组织为由局部 anchor、切空间和残差尺度构成的正常变化流形。
局部 anchor 表示当前图像中的正常外观模式，切空间表示这些模式附近允许的变化方向，残差尺度度量局部证据越过正常边界的程度。
每个目标 patch 随后被分解为流形内投影与流形外残差：流形内投影参与估计当前实例的 normality boundary，流形外残差被保留为异常定位证据。
整个过程在单张图像内完成，不使用目标训练图像、标签、mask、外部正常样本库或目标类别统计。

本文贡献如下：
\begin{enumerate}
  \item 指出 CLIP-ZSAD 中语义迁移与边界迁移的分离：现有方法主要提升 abnormal semantic transfer，而异常定位仍需要 instance-conditioned normality boundary。
  \item 将零样本异常定位重构为实例条件化正常性边界估计，把 normal 建模为当前实例中可接受局部变化的边界结构。
  \item 提出 \method{}，以正常变化流形表示该边界，将局部正常模式、允许变化方向和越界残差统一到一个机制中。
  \item 设计边界内投影与边界外残差分解，使 normal 证据能够适配当前实例，同时避免缺陷区域被错误吸收到 normal 表示中。
  \item 在 MVTec AD 上验证 \method{} 相比 AnomalyCLIP 的平均 P-AUROC、P-AUPRO、I-AUROC 和 I-AP 均有提升，并通过结构消融分析性能来源。
\end{enumerate}

\section{相关工作}
\label{sec:related_work}

\subsection{CLIP 零样本异常检测}

视觉语言模型为零样本异常检测提供了可迁移的状态语义表征。
WinCLIP 将状态词组合与窗口级局部特征用于零样本和少样本异常分类与分割 \cite{jeong2023winclip}。
APRIL-GAN 改善 \clip{} patch 与文本状态之间的局部对齐 \cite{chen2023aprilgan}。
AnomalyCLIP、AdaCLIP、VCP-CLIP 和 AA-CLIP 等方法进一步通过对象无关提示、混合提示、视觉上下文提示或轻量适配器增强 normal/abnormal 状态建模 \cite{zhou2024anomalyclip,yao2024adaclip,qu2024vcpclip,ma2025aaclip}。
这些工作共同回答了一个关键问题：如何让 abnormal 语义脱离具体目标类别，并在未见类别上保持可用。
\method{} 研究同一框架中更靠近定位决策的一环：当 abnormal 语义已经能够迁移时，normal 语义在当前实例中的可接受边界如何确定。

\subsection{正常性结构建模}

工业异常检测长期依赖正常样本的局部表示结构。
PaDiM 使用 patch 表示的分布统计刻画正常性 \cite{defard2021padim}，PatchCore 通过正常 patch 记忆库与近邻检索获得强定位性能 \cite{roth2022patchcore}，UniAD 探索多类别统一建模并分析重建捷径对异常定位的影响 \cite{you2022uniad}。
这些方法表明，正常性不能完全还原为一个二分类语义标签；它还包含局部邻域关系、分布尺度和残差解释能力。
然而，这类结构通常来自目标类别的正常训练样本、记忆库或统一训练过程。
\method{} 将结构化正常性的思想放入目标类别零样本协议：方法不访问目标训练图像，也不为目标类别维护正常样本库，而是在单张测试图像内部估计可被 normal 语义接受的局部变化结构。

\subsection{实例条件化正常性边界}

异常定位中的实例条件化首先是一个局部证据选择问题。
当前图像同时包含可接受正常变化和潜在缺陷，关键不在于让模型贴合整张图像，而在于选择哪些局部证据可以定义 normality boundary，哪些局部偏离必须作为异常证据保留。
因此，直接聚合低异常响应区域会带来一个结构性风险：如果缺陷区域被吸收到实例 normal 表示中，定位信号就会被边界更新抹平。
\method{} 将实例条件化组织为边界估计问题。
流形内投影只吸收与当前正常变化相容的证据，流形外残差则显式保留不能由该边界解释的局部偏离。
这一设计使 normal 语义获得实例尺度的边界信息，同时避免把异常证据转化为新的正常原型。

\section{方法}
\label{sec:method}

\method{} 的目标是在不访问目标类别正常样本库的条件下，为每张测试图像估计一个临时的正常性边界。
该边界由当前图像中被 CLIP normal 语义支持、并且在空间上保持一致的局部证据定义。
方法首先用冻结的 \clip{} 提供 normal/abnormal 状态语义，再从当前图像中构建正常变化流形。
流形内方向定义可被当前实例接受的正常变化，流形外残差定义应被保留的异常证据。

\subsection{问题设置}
\label{sec:setting}

给定冻结的 \clip{} 图像编码器 \(f_\phi\)、文本编码器 \(g_\psi\)，以及 normal/abnormal 文本提示集合，目标是在未见目标类别上输出图像级异常分数和像素级异常图。
对于输入图像 \(x\)，从选定的视觉层提取 patch 表示并记为
\begin{equation}
  Z(x)=\{z_i\}_{i=1}^{N}, \qquad z_i \in \R^d,\qquad \norm{z_i}_2=1 .
  \label{eq:patch_tokens}
\end{equation}
文本编码器产生 normal 与 abnormal 文本锚 \(t_n,t_a\in\R^d\)。
目标测试阶段只访问当前测试图像和文本提示，不使用目标训练图像、类别标签、异常标签、mask、外部正常样本库或累计测试流统计。
\method{} 为每张输入图像估计一次实例条件化正常性边界，该边界只服务于当前图像，推理完成后即被丢弃。

\subsection{单图像正常变化流形}
\label{sec:manifold}

\method{} 将当前实例的正常性边界表示为单图像正常变化流形。
这一表示的出发点是，正常性边界不能由单个 normal 原型充分刻画。
一个原型只能描述正常证据的中心，却不能说明该中心附近哪些变化方向仍然合法，也不能定义越界残差的尺度。
因此，正常变化流形由一组局部正常 anchor、对应的可容许变异方向和残差尺度组成。
局部 anchor 表示当前图像中的正常外观模式，切空间描述该模式附近的可容许变化，残差尺度用于规范化偏离该局部结构的流形外分量。
在这一表示下，每个 patch 都可以被解释为边界内投影与边界外残差的组合。
具体地，我们首先计算每个 patch 的语义异常先验：
\begin{equation}
  p_i^a=
  \softmax
  \left(
  \frac{1}{T}
  \begin{bmatrix}
  \langle z_i,t_n\rangle\\
  \langle z_i,t_a\rangle
  \end{bmatrix}
  \right)_2 ,
  \label{eq:semantic_prior}
\end{equation}
其中 \(T\) 是温度。
正常 anchor 由语义 normal 置信度和局部一致性共同确定。
令 \(\mathcal{N}(i)\) 表示 patch \(i\) 的空间邻域，定义
\begin{equation}
  h_i=
  \exp
  \left(
  -\frac{1}{|\mathcal{N}(i)|\sigma^2}
  \sum_{j\in\mathcal{N}(i)}
  \norm{z_i-z_j}_2^2
  \right),
  \qquad
  b_i=(1-p_i^a)^\eta h_i .
  \label{eq:normal_candidate}
\end{equation}
\(b_i\) 高的 patch 既接近 normal 文本锚，又与局部邻域保持一致，因此更适合作为当前图像中的正常边界证据。
从 \(b_i\) 最高的候选集合中选取 \(K\) 个局部 anchor，并在候选 patch 的近邻内估计切方向与残差尺度，得到单图像正常变化流形：
\begin{equation}
  \M=\{(a_k,U_k,s_k)\}_{k=1}^{K}.
  \label{eq:atlas}
\end{equation}
其中 \(a_k\in\R^d\) 是第 \(k\) 个局部正常 anchor，\(U_k\in\R^{d\times m}\) 是局部切空间基，\(s_k>0\) 是该邻域的残差尺度。
实现中，anchor 由当前图像的高 \(b_i\) 候选 patch 聚类或非极大值筛选得到；每个 anchor 周围的候选近邻 patch 经中心化后做低秩 SVD，前 \(m\) 个方向作为该邻域允许的正常变化方向。
局部尺度由切空间外残差的稳健分位数与邻域距离共同估计。
这些邻域与目标类别名称解耦，由当前图像内的高置信正常候选估计，用于刻画可由 normal 语义接受的边缘连续性、纹理方向、材质波动和成像变化。

给定测试 patch \(z_i\)，其正常邻域分配由最大余弦相似度确定：
\begin{equation}
  k_i=\arg\max_{k} \langle z_i,a_k\rangle .
  \label{eq:assignment}
\end{equation}
令 \(\Delta_i=z_i-a_{k_i}\)，则其流形内投影和流形外残差分别为
\begin{align}
  \hat z_i
  &=
  \frac{a_{k_i}+U_{k_i}U_{k_i}^{\top}\Delta_i}
       {\norm{a_{k_i}+U_{k_i}U_{k_i}^{\top}\Delta_i}_2}, \label{eq:projection}\\
  r_i
  &=
  \frac{\norm{\Delta_i-U_{k_i}U_{k_i}^{\top}\Delta_i}_2}
       {s_{k_i}+\epsilon}. \label{eq:residual}
\end{align}
残差 \(r_i\) 衡量当前 patch 中无法由局部正常变异解释的分量。
我们进一步定义流形相容度
\begin{equation}
  c_i=\exp(-r_i/\tau),
  \label{eq:compatibility}
\end{equation}
其中 \(\tau\) 是残差温度。
高 \(c_i\) 表示该 patch 与当前正常变异结构相容；高 \(r_i\) 表示该 patch 包含更强的边界外残差。

\subsection{实例条件化正常性边界估计}
\label{sec:anchor_calibration}

实例条件化正常性边界由 CLIP 状态语义和单图像正常变异结构共同确定。
\clip{} 的 normal/abnormal 锚提供开放词汇状态判断，流形内投影提供当前实例中可被 normal 边界接受的视觉证据。
\method{} 只聚合与正常变化流形相容的边界内证据，而不直接吸收高残差区域。
这样得到的实例 normal 表示描述的是当前图像中 normal 语义的可接受范围，而不是整张图像的平均外观。
基于式~\eqref{eq:semantic_prior} 的语义异常先验，用流形相容度和 normal 语义置信度共同定义边界证据权重：
\begin{equation}
  q_i=c_i(1-p_i^a)^\gamma .
  \label{eq:normal_weight}
\end{equation}
\(\gamma\) 控制语义 normal 置信度对边界证据的影响。
从 \(q_i\) 最高的 top-\(k\) patch 集合 \(\mathcal{S}\) 中聚合当前图像的实例正常锚：
\begin{equation}
  v_n=
  \frac{\sum_{i\in \mathcal{S}} q_i \hat z_i}
       {\norm{\sum_{i\in \mathcal{S}} q_i \hat z_i}_2}.
  \label{eq:visual_normal_anchor}
\end{equation}
\(v_n\) 由投影后的 \(\hat z_i\) 构造。
投影表征位于当前实例的正常变异邻域内，因此能够为边界估计提供实例化 normal 证据，并降低高残差局部区域对边界的影响。

设 \(\bar q\) 为 \(\mathcal{S}\) 内平均边界证据置信度，normal 语义锚按下式被实例条件化：
\begin{equation}
  t'_n=
  \frac{(1-\alpha \bar q)t_n+\alpha \bar q v_n}
       {\norm{(1-\alpha \bar q)t_n+\alpha \bar q v_n}_2},
  \qquad
  t'_a=t_a .
  \label{eq:text_anchor_update}
\end{equation}
其中 \(\alpha\) 是更新强度。
默认设置保持 abnormal 锚不变，因为异常形态在目标类别中具有开放性和稀疏性。
该估计不需要反向传播，也不更新视觉编码器。
它为每张图像生成临时边界锚 \((t'_n,t'_a)\)，使 normality boundary 对齐当前图像中被正常变化流形接受的局部证据。

\subsection{边界外残差定位}
\label{sec:scoring}

实例条件化后的语义异常分数为
\begin{equation}
  A_i^{\mathrm{sem}}=
  \softmax
  \left(
  \frac{1}{T}
  \begin{bmatrix}
  \langle z_i,t'_n\rangle\\
  \langle z_i,t'_a\rangle
  \end{bmatrix}
  \right)_2 .
  \label{eq:semantic_score}
\end{equation}
流形外残差分数由 \(r_i\) 经图像内归一化得到：
\begin{equation}
  A_i^{\mathrm{res}}=\Norm(r_i).
  \label{eq:residual_score}
\end{equation}
最终 patch 级异常分数融合两类证据：
\begin{equation}
  A_i=(1-\rho)A_i^{\mathrm{sem}}+\rho A_i^{\mathrm{res}},
  \label{eq:final_score}
\end{equation}
其中 \(\rho\) 是残差权重。
语义分数继承 \clip{} 的状态判别能力，流形残差提供正常变异无法解释的局部证据。
像素级 anomaly map 由 patch 分数双线性上采样得到，图像级异常分数由 top-\(q\) patch 分数均值给出：
\begin{equation}
  A_{\mathrm{img}}=
  \frac{1}{q}
  \sum_{i\in \topk(A,q)} A_i .
  \label{eq:image_score}
\end{equation}

\begin{figure}[H]
  \centering
  \fbox{
  \begin{minipage}{0.91\linewidth}
  \small
  \textbf{冻结编码：}
  测试图像 \(x\) \(\rightarrow\) \clip{} image encoder \(\rightarrow\) patch tokens；
  normal/abnormal prompts \(\rightarrow\) \clip{} text encoder \(\rightarrow\) 文本锚。

  \vspace{1.5mm}
  \textbf{单图像推理：}
  文本引导正常候选 \(\rightarrow\) 临时正常变化流形 \(\rightarrow\) 流形投影与残差分解。
  流形内投影聚合为实例正常锚并估计当前图像的 normality boundary；流形外残差与实例条件化语义分数融合为 anomaly map。
  \end{minipage}
  }
  \caption{\method{} 方法总览。正常变化流形同时刻画当前实例的边界内正常证据和边界外异常残差。}
  \label{fig:overview}
\end{figure}

\begin{table}[H]
  \centering
  \caption{\method{} 的单图像推理流程。实例条件化边界只对当前输入图像有效。}
  \label{tab:algorithm}
  \small
  \begin{tabularx}{\linewidth}{p{0.16\linewidth}Y}
    \toprule
    步骤 & 操作 \\
    \midrule
    输入 & 测试图像 \(x\)，冻结 \clip{}，normal/abnormal 文本提示 \\
    单图像图谱 & 由 \clip{} image encoder 得到 patch tokens，由 text encoder 得到文本锚，并用高 normal 置信度、局部一致的 patch 构建正常变化流形 \(\M\) \\
    流形分解 & 对每个 patch 计算邻域分配、流形内投影 \(\hat z_i\)、流形外残差 \(r_i\) 和相容度 \(c_i\) \\
    边界证据 & 用 \(c_i(1-p_i^a)^\gamma\) 选择高相容 normal patch，并聚合投影后的 \(\hat z_i\) 得到实例正常锚 \(v_n\) \\
    边界估计 & 将 \(t_n\) 与 \(v_n\) 做置信度加权融合，得到当前图像的临时边界锚 \(t'_n\) \\
    异常评分 & 融合实例条件化语义异常分数与流形外残差分数，输出 anomaly map 和图像级异常分数 \\
    \bottomrule
  \end{tabularx}
\end{table}

\section{实验}
\label{sec:experiments}

\subsection{实验设置}
\label{sec:protocol}

我们在 MVTec AD \cite{bergmann2019mvtec} 上评估 \method{}。
实验采用严格的目标类别零样本协议：目标测试类别不提供训练图像、类别标签、异常标签或 mask；方法不使用外部正常样本库，也不从测试流累计目标类别统计。
每张图像独立构建正常变化流形，推理结束后丢弃。
所有结果使用相同的 \clip{} backbone、输入分辨率和 normal/abnormal 文本模板。

图像级指标包括 I-AUROC 和 I-AP；像素级指标包括 P-AUROC 和 P-AUPRO。
P-AUROC 衡量像素级排序质量，P-AUPRO 更强调缺陷区域覆盖与正常区域误报之间的平衡。
本文以 AnomalyCLIP 作为主要 CLIP-ZSAD 参照，并报告 \method{} 在 15 个类别上的完整结果。

\subsection{主结果}
\label{sec:main_results}

表~\ref{tab:mvtec_results} 给出 MVTec AD 上的类别级结果，并报告相对 AnomalyCLIP 的逐类差值。
\method{} 将 AnomalyCLIP 的平均 P-AUROC/P-AUPRO/I-AUROC/I-AP 从 91.1/81.4/91.6/96.4 提升到 92.9/84.6/92.4/96.7。
增益主要体现在像素级定位指标上，尤其是 cable、transistor、metal\_nut 和 zipper 等正常外观变异较大或缺陷形态较细碎的类别。
这表明实例条件化正常性边界能够更准确地区分可容许正常变异与边界外缺陷残差。

\begin{table}[H]
  \centering
  \caption{MVTec AD 上的 \method{} 类别级结果。 \(\Delta\) 表示相对 AnomalyCLIP 的差值，所有指标越高越好。}
  \label{tab:mvtec_results}
  \scriptsize
  \setlength{\tabcolsep}{3.2pt}
  \begin{tabular}{lrrrrrrrr}
    \toprule
    \multirow{2}{*}{类别} &
    \multicolumn{2}{c}{P-AUROC} &
    \multicolumn{2}{c}{P-AUPRO} &
    \multicolumn{2}{c}{I-AUROC} &
    \multicolumn{2}{c}{I-AP} \\
    \cmidrule(lr){2-3}\cmidrule(lr){4-5}\cmidrule(lr){6-7}\cmidrule(lr){8-9}
    & \method{} & \(\Delta\) & \method{} & \(\Delta\) & \method{} & \(\Delta\) & \method{} & \(\Delta\) \\
    \midrule
    carpet & 99.0 & +0.2 & 91.9 & +1.9 & 100.0 & +0.0 & 100.0 & +0.0 \\
    bottle & 92.3 & +1.9 & 84.2 & +3.4 & 90.4 & +1.7 & 97.3 & +0.5 \\
    hazelnut & 97.7 & +0.5 & 93.5 & +1.0 & 97.6 & +0.4 & 98.7 & +0.2 \\
    leather & 98.8 & +0.2 & 93.6 & +1.4 & 99.8 & +0.0 & 99.9 & +0.0 \\
    cable & 83.6 & +4.7 & 70.7 & +6.7 & 74.5 & +4.2 & 84.3 & +2.6 \\
    capsule & 96.4 & +0.6 & 89.6 & +2.0 & 89.5 & +0.0 & 97.8 & +0.0 \\
    grid & 97.6 & +0.3 & 80.1 & +4.7 & 97.9 & +0.1 & 99.3 & +0.0 \\
    pill & 93.4 & +1.6 & 90.3 & +2.2 & 83.3 & +2.2 & 95.9 & +0.6 \\
    transistor & 76.8 & +6.0 & 64.1 & +5.9 & 94.2 & +0.3 & 92.6 & +0.5 \\
    metal\_nut & 80.5 & +5.9 & 77.3 & +6.2 & 93.5 & +1.1 & 98.5 & +0.3 \\
    screw & 97.8 & +0.3 & 89.3 & +1.3 & 82.1 & +0.0 & 92.9 & +0.0 \\
    toothbrush & 93.8 & +1.9 & 89.9 & +1.4 & 87.5 & +2.2 & 94.8 & +0.9 \\
    zipper & 93.0 & +1.7 & 72.5 & +7.1 & 98.5 & +0.1 & 99.5 & +0.0 \\
    tile & 95.5 & +0.8 & 89.0 & +1.6 & 100.0 & +0.0 & 100.0 & +0.0 \\
    wood & 96.6 & +0.2 & 92.4 & +0.9 & 97.3 & +0.4 & 99.3 & +0.1 \\
    \midrule
    mean & \textbf{92.9} & \textbf{+1.8} & \textbf{84.6} & \textbf{+3.2} & \textbf{92.4} & \textbf{+0.8} & \textbf{96.7} & \textbf{+0.3} \\
    \bottomrule
  \end{tabular}
\end{table}

\subsection{消融分析}
\label{sec:ablation}

表~\ref{tab:ablation} 分析实例条件化正常性边界的结构贡献。
固定边界对应 AnomalyCLIP 的原始 normal/abnormal 语义判断。
单原型边界用当前图像低异常响应区域构造一个实例 normal 锚，主要调整边界中心。
多锚点边界进一步覆盖实例内部的多种正常外观模式；引入局部切空间后，模型获得边界内变异与边界外残差的显式分解。
相容权重控制哪些局部证据参与边界估计，边界外残差提供独立的定位信号。
完整 \method{} 同时保留局部 anchor、切空间投影、相容权重和边界外残差，因此在四个指标上取得最高平均结果。

\begin{table}[H]
  \centering
  \caption{MVTec AD 核心消融。所有指标越高越好。}
  \label{tab:ablation}
  \small
  \begin{tabular}{lcccc}
    \toprule
    变体 & P-AUROC & P-AUPRO & I-AUROC & I-AP \\
    \midrule
    固定 normal/abnormal 边界 & 91.1 & 81.4 & 91.6 & 96.4 \\
    单锚点实例边界 & 91.5 & 80.8 & 92.0 & 96.5 \\
    多锚点实例边界 & 92.0 & 82.7 & 92.1 & 96.5 \\
    \method{} 去掉相容权重 & 91.8 & 82.1 & 92.0 & 96.5 \\
    \method{} 去掉边界外残差 & 92.3 & 83.0 & 92.3 & 96.6 \\
    完整 \method{} & \textbf{92.9} & \textbf{84.6} & \textbf{92.4} & \textbf{96.7} \\
    \bottomrule
  \end{tabular}
\end{table}

\subsection{参数分析}
\label{sec:sensitivity}

表~\ref{tab:sensitivity} 展示单图像正常变化流形规模与切空间秩的影响。
较小的 \(K\) 会欠拟合实例内的正常外观模式，降低残差响应的空间分辨率；过大的 \(K\) 会增加检索开销并放大局部尺度估计噪声。
零秩切空间退化为局部原型匹配，难以刻画纹理方向和材质连续变化。
在当前设置下，\(K=64\)、\(m=8\) 在边界精度与推理开销之间取得较好平衡。

\begin{table}[H]
  \centering
  \caption{单图像正常变化流形的参数敏感性。时间为单张图像额外推理开销。}
  \label{tab:sensitivity}
  \small
  \begin{tabular}{lccc}
    \toprule
    设置 & P-AUROC & P-AUPRO & 额外时间 \\
    \midrule
    \(K=16,m=8\) & 92.0 & 82.7 & 11 ms \\
    \(K=32,m=8\) & 92.5 & 83.8 & 16 ms \\
    \(K=64,m=8\) & \textbf{92.9} & \textbf{84.6} & 23 ms \\
    \(K=128,m=8\) & 92.8 & 84.4 & 39 ms \\
    \(K=64,m=0\) & 92.1 & 82.9 & 18 ms \\
    \(K=64,m=4\) & 92.6 & 84.0 & 21 ms \\
    \(K=64,m=16\) & 92.8 & 84.3 & 29 ms \\
    \bottomrule
  \end{tabular}
\end{table}

\method{} 不进行反向传播，不更新视觉编码器，也不保存目标类别正常样本库。
额外计算主要来自 anchor 检索、局部投影和残差图上采样，可作为 CLIP-ZSAD 推理过程中的轻量边界估计步骤。

\subsection{定性分析}
\label{sec:qualitative}

图~\ref{fig:qualitative} 概括典型定位现象。
固定边界容易在颜色、光照或纹理方向偏移的正常区域产生零散误报，并在低对比缺陷上形成分散响应。
\method{} 将局部证据分解为可容许正常变异和边界外残差，使正常区域响应被抑制，同时在划痕、孔洞边界和局部污染区域形成更集中的热图。
这种差异在小缺陷和边界缺陷上尤其明显，因为这些区域占比低，直接聚合全图证据时容易被大面积正常区域稀释。

\begin{figure}[H]
  \centering
  \fbox{
  \begin{minipage}{0.91\linewidth}
  \small
  \begin{tabularx}{\linewidth}{>{\bfseries}p{0.18\linewidth}YYYY}
    场景 & 固定边界 & 直接边界更新 & \method{} & 机制响应 \\
    \midrule
    正常成像变化 & 背景纹理与光照区域出现零散误报 & 误报减少但缺陷边界变钝 & 正常区域响应被压低 & 高相容权重覆盖正常纹理变化 \\
    低对比小缺陷 & 缺陷响应弱且热图分散 & 缺陷局部容易进入 normal 锚 & 缺陷区域形成集中响应 & 边界外残差增强缺陷边界 \\
  \end{tabularx}
  \end{minipage}
  }
  \caption{定性分析示意。\method{} 同时呈现正常变化流形接受区域与流形外残差响应区域。}
  \label{fig:qualitative}
\end{figure}

\section{讨论}
\label{sec:discussion}

\method{} 将 CLIP-ZSAD 中的状态语义比较扩展为实例条件化正常性边界估计。
固定状态比较以同一组 normal/abnormal 锚评估所有未见实例，适合表达对象无关异常语义，却缺少当前实例中 normal 可接受范围的信息。
\method{} 在当前图像中估计正常变化流形，并据此选择能够支撑 normality boundary 的局部证据。
这种结构使 normal 语义能够吸收当前实例的合法外观变化，同时将无法由正常变化解释的残差保留为定位信号。
实验增益主要体现在 P-AUPRO 和困难类别定位质量上，与边界外残差建模的设计目标一致。

当前方法依赖单张图像中存在足够的可靠正常区域来估计正常变化流形。
当缺陷覆盖范围极大、正常区域极少，或同一图像包含多个强烈冲突的正常外观模式时，边界估计会受到候选 patch 质量限制。
这一限制也刻画了本文主张的适用边界：NCMA 估计的是实例内部可观察正常变化所诱导的边界，而不是从外部正常样本库恢复完整类别分布。
面向视频、批量检测或多视角工业场景时，可以进一步把单图像流形扩展为时序一致或跨视角一致的正常性边界，同时保持不使用目标缺陷标签的零样本协议。

\section{结论}
\label{sec:conclusion}

本文将 CLIP 零样本异常检测中的固定状态比较重构为实例条件化正常性边界估计。
这一重构区分了 abnormal 语义的跨类别迁移和 normality boundary 的实例化估计。
\method{} 在单张图像内构建文本引导的正常变化流形，并将局部 patch 分解为边界内投影和边界外残差。
边界内投影用于估计当前实例中 normal 语义的可接受范围，边界外残差与实例条件化语义分数共同形成 anomaly map。
在 MVTec AD 上，\method{} 将 AnomalyCLIP 的平均 P-AUROC/P-AUPRO/I-AUROC/I-AP 从 91.1/81.4/91.6/96.4 提升到 92.9/84.6/92.4/96.7。
这些结果表明，正常变化流形为 CLIP-ZSAD 提供了实例边界机制，能够建模合法正常变化并保留局部缺陷证据。

\begin{thebibliography}{99}

\bibitem{radford2021clip}
Alec Radford, Jong Wook Kim, Chris Hallacy, Aditya Ramesh, Gabriel Goh, Sandhini Agarwal, Girish Sastry, Amanda Askell, Pamela Mishkin, Jack Clark, Gretchen Krueger, and Ilya Sutskever.
\newblock Learning transferable visual models from natural language supervision.
\newblock In \emph{ICML}, 2021.

\bibitem{jeong2023winclip}
Jongheon Jeong, Yang Zou, Taewan Kim, Dongqing Zhang, Avinash Ravichandran, and Onkar Dabeer.
\newblock WinCLIP: Zero-/few-shot anomaly classification and segmentation.
\newblock In \emph{CVPR}, 2023.

\bibitem{chen2023aprilgan}
Xuhai Chen, Yue Han, and Jiangning Zhang.
\newblock APRIL-GAN: A zero-/few-shot anomaly classification and segmentation method for CVPR 2023 VAND workshop challenge.
\newblock \emph{arXiv preprint arXiv:2305.17382}, 2023.

\bibitem{zhou2024anomalyclip}
Qihang Zhou, Guansong Pang, Yu Tian, Shibo He, and Jiming Chen.
\newblock AnomalyCLIP: Object-agnostic prompt learning for zero-shot anomaly detection.
\newblock \emph{arXiv preprint arXiv:2310.18961}, 2024.

\bibitem{yao2024adaclip}
Xincheng Yao, Ruoqi Li, Zefeng Qian, Lu Wang, and Chongyang Zhang.
\newblock AdaCLIP: Adapting CLIP with hybrid learnable prompts for zero-shot anomaly detection.
\newblock \emph{arXiv preprint arXiv:2407.15795}, 2024.

\bibitem{qu2024vcpclip}
Xiaozhen Qu, Ming Gao, and colleagues.
\newblock VCP-CLIP: A visual context prompting model for zero-shot anomaly segmentation.
\newblock \emph{arXiv preprint arXiv:2407.12276}, 2024.

\bibitem{ma2025aaclip}
Xinyu Ma, and colleagues.
\newblock AA-CLIP: Enhancing zero-shot anomaly detection via anomaly-aware CLIP.
\newblock \emph{arXiv preprint arXiv:2503.06661}, 2025.

\bibitem{defard2021padim}
Thomas Defard, Aleksandr Setkov, Angelique Loesch, and Romaric Audigier.
\newblock PaDiM: A patch distribution modeling framework for anomaly detection and localization.
\newblock In \emph{ICPR}, 2021.

\bibitem{roth2022patchcore}
Karsten Roth, Latha Pemula, Joaquin Zepeda, Bernhard Schölkopf, Thomas Brox, and Peter Gehler.
\newblock Towards total recall in industrial anomaly detection.
\newblock In \emph{CVPR}, 2022.

\bibitem{you2022uniad}
Zhiyuan You, Kai Yang, Wenhan Luo, Lei Cui, Yujun Zheng, and Xiu Li.
\newblock A unified model for multi-class anomaly detection.
\newblock In \emph{NeurIPS}, 2022.

\bibitem{bergmann2019mvtec}
Paul Bergmann, Michael Fauser, David Sattlegger, and Carsten Steger.
\newblock MVTec AD: A comprehensive real-world dataset for unsupervised anomaly detection.
\newblock In \emph{CVPR}, 2019.

\end{thebibliography}

\end{document}
